\chapter*{Puente de protocolización hacia LAB-FEN-001}
\addcontentsline{toc}{chapter}{Puente de protocolización hacia LAB-FEN-001}

\section*{De la preparación al ensayo interno}

\begin{observation}
Los criterios de protocolización y trazabilidad desarrollados en este tomo pueden aplicarse a una pieza concreta sin confundir esa aplicación con una propiedad universal del canon. LAB-FEN-001 constituye una aplicación interna de este principio.
\end{observation}

\begin{observation}
La batería distingue controles denominados \texttt{LAB-D0}, \texttt{LAB-D1} y \texttt{LAB-D2}. Estas etiquetas pertenecen exclusivamente al laboratorio y no sustituyen ni reinterpretan los niveles dimensionales D0--D11.
\end{observation}

\section*{Paso al gate metodológico}

\begin{observation}
Bajo el contrato congelado de LAB-FEN-001-v1, \texttt{LAB-D2-v1} superó su firma local, las ablaciones primarias predeclaradas y el atlas paramétrico local. Posteriormente fue sometido al gate DMO-004.
\end{observation}

\begin{observation}
Los siete criterios de DMO-004 resultaron satisfechos. La promoción máxima permitida por ese procedimiento es \texttt{REGIMEN INTERNO ADMITIDO BAJO LAB-FEN-001-v1 Y DMO-004}. Esta declaración es metodológica e interna y no equivale a validación física externa.
\end{observation}

\section*{Función de Tomo III}

\begin{observation}
La caracterización cuantitativa, las ablaciones, el atlas, la profundidad histórica y los límites de interpretación se concentran en Tomo III. El presente tomo sólo fija el puente entre preparación operativa y tratamiento metodológico auditado.
\end{observation}
